% ============================================================================
% Fichier  : partie4/P04_C11_smsi_mettre_oeuvre_auditer.tex
% Titre    : Mettre en Œuvre et Auditer un SMSI
% Auteur   : MINKA MI NGUIDJOÏ Thierry Emmanuel
% Version  : 1.0 -- novembre 2026
% Corrigés : corriges/Corr_C11.tex
% ============================================================================
% Sources : ISO/IEC 27001:2022 (art. 4.3, 6.2, 7.5, 9.1 à 9.3, 10.2 ;
% informations documentées exigées ; éléments d'entrée de la revue de
% direction : paraphrasés) ; ISO/IEC 27004:2016 (mesures de mise en œuvre et
% d'efficacité) ; ISO/IEC 27007:2020 ; ISO/IEC 17021-1:2015 (non-conformités
% majeure et mineure, définitions paraphrasées). Loi n° 2010/012, art. 26 et
% 30 (résumés, Chapitre 10).
% Faits du cas (bible) : BANQUE SANAGA, extension du SMSI aux agences,
% renouvellement du certificat en juin 2028, audit interne de mars 2028,
% déclaration d'applicabilité corrigée (C08), R-AG-02 (GAB) à apprécier.
% Faits nouveaux (reportés dans la bible) : projet d'extension de mars 2027
% à mai 2028 ; périmètre 2028 (monétique exclue, interfaces incluses) ; six
% indicateurs I1 à I6 (base 4/1/1) ; audit interne du SMSI de mai 2028 ;
% revue de direction du 30 mai 2028 ; étape 2 de CERTIQUAL du 12 au
% 16 juin 2028 : 1 non-conformité majeure (R-AG-02 non apprécié), 3
% mineures, 2 pistes de progrès ; correction acceptée le 25 août 2028 ;
% certificat renouvelé en septembre 2028, avec un trou de trois mois.
% ============================================================================

\chapter{Mettre en Œuvre et Auditer un SMSI}
\label{chap:smsi}
\chaptermark{Mettre en œuvre et auditer un SMSI}

\epigraph{%
    « Un système de management qui ne se mesure pas\\
    ne se gère pas~; qui ne s'audite pas ne se croit pas. »%
}{Principe du cours}

\niveauChapitre{Intermédiaire}{%
    Maîtrise du Chapitre~\ref{chap:iso27000}, du Chapitre~\ref{chap:conclure}%
}

\begin{boxobjectifs}
\begin{itemize}
    \item Conduire la mise en œuvre d'un SMSI comme un projet, de la
          décision de la direction à l'audit de certification.
    \item Définir un périmètre défendable et en justifier les exclusions.
    \item Identifier les informations documentées qu'exige
          l'ISO/IEC~27001.
    \item Construire des indicateurs selon l'ISO/IEC~27004, et coter leur
          fiabilité.
    \item Conduire l'audit interne du SMSI et préparer la revue de
          direction.
    \item Qualifier une non-conformité majeure ou mineure, et en tirer les
          conséquences pour la certification.
\end{itemize}
\end{boxobjectifs}

\bigskip

% ============================================================================
\section*{Avant de commencer : « Majeure ou mineure~? »}
% ============================================================================

Douala, 16~juin 2028, réunion de clôture de l'étape~2 de l'audit de
renouvellement\footnote{Comme dans tout le cours, la BANQUE SANAGA,
CERTIQUAL AFRIQUE, les dates et les montants sont fictifs~; les normes
citées sont exactes et paraphrasées.}. Le chef de mission de CERTIQUAL
AFRIQUE lit ses constats. Au troisième, la RSSI se raidit~: « La banque
n'a pas apprécié le risque des automates de retrait de quatorze agences,
alors que l'audit interne l'avait signalé en avril. Je qualifie cette
non-conformité de \emph{majeure}. »

Le directeur général demande ce que cela change. Tout~: le certificat
de 2025 expire dans dix jours, et l'organisme ne peut pas le renouveler
tant que la non-conformité majeure n'est pas corrigée et sa correction
vérifiée. Pendant ce temps, la banque n'aura plus de certificat.

\medskip

Ce chapitre ouvre la Partie~IV, consacrée à la pratique. Il suit
l'extension du SMSI de la BANQUE SANAGA à toute la banque, de la décision du
conseil à l'audit de certification, et explique pourquoi un seul constat a
suffi à interrompre la certification.

\begin{boxalert}
\textbf{Violation typique.} Traiter l'audit de certification comme un
examen à préparer, en mettant à jour les documents dans les semaines qui
précèdent, au lieu de faire fonctionner le système toute l'année.

\smallskip
\textbf{Règle~:} l'organisme certificateur n'évalue pas des documents~; il
évalue un système qui fonctionne, par ses traces. Une exigence qui n'a pas
de trace sur la période n'est pas satisfaite, même si le document existe.
\end{boxalert}

\separateur

% ============================================================================
\section{Mettre en Œuvre : un Projet, puis un Système}
\label{sec:c11-projet}
% ============================================================================

\subsection{Les étapes}

L'ISO/IEC~27001 décrit des exigences, pas un projet. Dans la pratique, la
mise en œuvre d'un SMSI suit pourtant un ordre presque toujours le même,
parce que chaque étape a besoin de la précédente.

\begin{tableaucours}{cL{4.4cm}L{3.2cm}L{3.6cm}}{%
    \tableauHeader{N°} &
    \tableauHeader{Étape} &
    \tableauHeader{Article} &
    \tableauHeader{BANQUE SANAGA}
}
1 & Décision de la direction, moyens & 5 & Conseil, octobre 2027 \\
2 & Contexte, parties intéressées, périmètre & 4 & Décembre 2026, revu en
2027 \\
3 & Politique et objectifs de sécurité & 5.2, 6.2 & Politique révisée
novembre 2027 \\
4 & Appréciation et traitement des risques, déclaration d'applicabilité &
6.1, 8.2, 8.3 & Décembre 2026 à avril 2028 \\
5 & Compétences, sensibilisation, documentation & 7 & 2027--2028 \\
6 & Exploitation des mesures & 8.1 & 2027--2028 \\
7 & Surveillance, audit interne, revue de direction & 9 & Mars à mai
2028 \\
8 & Non-conformités et améliorations & 10 & Continu \\
\end{tableaucours}
\vspace{-0.8em}
\begin{center}{\small\itshape Tableau~11.1. Les étapes de la mise en œuvre
d'un SMSI}\end{center}

Une fois les étapes franchies, le projet devient un système~: les
étapes~4 à~8 se répètent chaque année, et à chaque changement important.
C'est ce que l'organisme certificateur vient vérifier.

\subsection{Un périmètre défendable}

Le prologue a montré un périmètre de complaisance. L'ISO/IEC~27001 demande
de fixer le périmètre en tenant compte des enjeux, des exigences des
parties intéressées, et des interfaces et dépendances avec ce qui reste en
dehors~\cite{iso27001_2022}.

\begin{boxscenario}{Opération SANAGA. Le périmètre de 2028}
\textbf{Rédaction approuvée par le conseil.} «~Le système de management de
la sécurité de l'information couvre les activités bancaires de la BANQUE
SANAGA~: siège, centre de données principal de Douala, site de secours de
Bonabéri, et 38 agences, avec leurs postes, leurs réseaux, la messagerie et
les accès des prestataires. L'activité monétique est exclue~: elle relève
d'un programme de conformité PCI DSS distinct~; ses interfaces avec le
système central sont incluses.~»

\medskip
\textbf{Analyse} (confiance élevée). L'exclusion de la monétique est
défendable pour trois raisons~: elle est couverte par un autre dispositif,
nommé~; les interfaces, par lesquelles une attaque passerait de l'une à
l'autre, sont dans le périmètre~; et la décision est approuvée par
l'instance qui répond des risques. C'est l'inverse de 2025, où les agences
avaient été exclues sans dispositif de remplacement, sans maîtrise des
interfaces et sans décision du conseil.
\end{boxscenario}

\subsection{Les informations documentées}

\begin{boxprocedure}
\textbf{Ce que la norme exige d'écrire.} Le dossier du SMSI contient au
moins les informations documentées que l'ISO/IEC~27001 exige
expressément~:
\begin{itemize}
    \item le périmètre (4.3) et la politique de sécurité (5.2)~;
    \item le processus d'appréciation des risques (6.1.2) et ses résultats
          (8.2)~; le processus de traitement (6.1.3), la déclaration
          d'applicabilité et le plan de traitement, et ses résultats (8.3)~;
    \item les objectifs de sécurité (6.2)~;
    \item les preuves de compétence (7.2)~;
    \item les informations nécessaires à la maîtrise des processus (8.1)~;
    \item les résultats de la surveillance et de la mesure (9.1)~;
    \item le programme et les résultats des audits internes (9.2)~;
    \item les résultats des revues de direction (9.3)~;
    \item les non-conformités, les actions menées et leurs résultats
          (10.2).
\end{itemize}
Tout le reste est laissé au jugement de l'organisation (7.5)~: écrire ce
qui est nécessaire à l'efficacité, pas davantage.
\end{boxprocedure}

% ============================================================================
\section{Mesurer : les Indicateurs du SMSI}
\label{sec:c11-indicateurs}
% ============================================================================

\subsection{Deux familles de mesures}

L'article~9.1 demande de déterminer ce qui doit être surveillé et mesuré,
comment, quand, par qui, et comment les résultats sont analysés.
L'ISO/IEC~27004 aide à construire ces mesures~\cite{iso27004_2016}. Elle
distingue deux familles~:
\begin{itemize}
    \item les \textbf{mesures de mise en œuvre}~: la mesure est-elle en
          place, et à quel degré~? (part des comptes avec second facteur)~;
    \item les \textbf{mesures d'efficacité}~: la mesure produit-elle
          l'effet attendu~? (part des agents qui cliquent lors d'une
          campagne d'hameçonnage simulé).
\end{itemize}

Un tableau de bord fait seulement de mesures de mise en œuvre dit que l'on
travaille~; il ne dit pas si l'on est protégé.

\subsection{Les six indicateurs de la RSSI}

\begin{tableaucours}{cL{3.8cm}cL{2.8cm}cc}{%
    \tableauHeader{Réf.} &
    \tableauHeader{Indicateur} &
    \tableauHeader{Famille} &
    \tableauHeader{Source} &
    \tableauHeader{Cible} &
    \tableauHeader{Mai 2028}
}
I1 & Comptes avec second facteur & M.~œ. & Extraction de l'annuaire &
100~\% & 100~\% \\
I2 & Délai médian d'application des correctifs critiques & Eff. & Outil
de déploiement & 15~j & 9~j \\
I3 & Restaurations trimestrielles réussies & Eff. & Procès-verbaux & 100~\%
& 100~\% \\
I4 & Agences qui envoient des journaux & M.~œ. & Collecteur, quotidien &
100~\% & 100~\% \\
I5 & Agences ayant exercé le mode dégradé dans l'année & M.~œ. &
Déclarations des agences & 100~\% & 79~\% \\
I6 & Agents qui cliquent lors d'une campagne simulée & Eff. & Plateforme de
campagne & $<$~5~\% & 11~\% \\
\end{tableaucours}
\vspace{-0.8em}
\begin{center}{\small\itshape Tableau~11.2. Le tableau de bord du SMSI
(M.~œ.~: mise en œuvre~; Eff.~: efficacité)}\end{center}

\subsection{Un indicateur vaut ce que vaut sa source}

\begin{boxlecture}{La fiabilité du tableau de bord}
\textbf{Constatation.} Mme~EKAMBI cote la fiabilité de chaque indicateur,
c'est-à-dire la capacité de sa source à mesurer ce que dit sa définition~:
\begin{itemize}
    \item \OUI{}~: I1, I3, I4, I6 (sources automatiques ou procès-verbaux
          signés)~;
    \item \NON{}~: I2~: l'outil de déploiement ne couvre que les serveurs~;
          les postes des agences, mis à jour à part, n'entrent pas dans le
          calcul, alors que la définition parle de «~tous les systèmes~»~;
    \item \NEK{}~: I5~: les agences déclarent leurs exercices, sans compte
          rendu ni chronométrage (cf. Kribi, Chapitre~\ref{chap:conclure}).
\end{itemize}
4~\OUI{}, 1~\NON{}, 1~\NEK{}~: $\mu \approx 0{,}67$~; $U \approx 0{,}17$~;
\[
H = -\left(\tfrac{2}{3}\log_2\tfrac{2}{3} + 2 \times
  \tfrac{1}{6}\log_2\tfrac{1}{6}\right) \approx 0{,}390 + 0{,}862
  \approx 1{,}25~\text{bit}.
\]

\textbf{Analyse} (confiance élevée). Les deux indicateurs les plus
rassurants en apparence posent problème~: I2 est bon parce qu'il mesure un
sous-ensemble~; I5, à 79~\%, repose sur des déclarations. Un indicateur
dont la source est déclarative est un \NEK{}, quelle que soit sa valeur. Le
plus utile du tableau est I6, à 11~\%~: c'est une mauvaise nouvelle, mais
une mesure fiable d'efficacité.

\medskip
\hfill\lecture{0{,}67}{0{,}17}{1{,}25}
\end{boxlecture}

\begin{boxinfo}
\textbf{La fiche d'un indicateur.} Chaque indicateur se décrit par~:
l'objectif de sécurité qu'il sert~; sa définition et sa formule~; sa source
et son mode de collecte~; sa fréquence~; sa cible et ses seuils d'alerte~;
son responsable~; et la cote de fiabilité de sa source. C'est la fiche
d'assertion (Chapitre~\ref{chap:conclure}) appliquée à une mesure.
\end{boxinfo}

% ============================================================================
\section{Auditer le SMSI}
\label{sec:c11-audit}
% ============================================================================

\subsection{Deux objets}

L'audit interne exigé par l'article~9.2 porte sur deux objets~: le
\emph{système de management} (les articles~4 à~10~: le périmètre est-il
tenu, les risques appréciés, les objectifs suivis, les non-conformités
traitées~?) et les \emph{mesures} de la déclaration d'applicabilité.
L'ISO/IEC~27007 donne les lignes directrices de cet audit~\cite{iso27007_2020}.
La mission de mars~2028 sur les agences (Chapitre~\ref{chap:audit}) portait
sur des mesures~; il manquait l'audit du système lui-même.

\begin{boxscenario}{Opération SANAGA. L'audit interne du SMSI, mai 2028}
\textbf{Programme.} L'audit interne examine, pour chaque article de 4 à 10,
une ou deux exigences, avec la preuve de leur fonctionnement sur la
période~: approbation du périmètre (4.3), revue de la politique (5.2),
appréciation des risques à jour (6.1.2, 8.2), suivi des objectifs (6.2),
preuves de compétence (7.2), indicateurs (9.1), traitement des
non-conformités de l'audit de mars (10.2).

\medskip
\textbf{Constat principal.} L'appréciation des risques n'a pas été mise à
jour après l'audit de mars~: le risque R-AG-02 (compromission d'un GAB),
signalé le 20~avril, n'a ni cote, ni propriétaire, ni décision. La RSSI
l'attendait «~pour la prochaine revue annuelle~».

\medskip
\textbf{Analyse} (confiance élevée). Le défaut est de même nature qu'en
2025~: une partie du périmètre comporte un risque connu qui n'a pas été
apprécié. Le rapport de l'audit interne le porte en tête, avec une
échéance avant l'étape~2.
\end{boxscenario}

\subsection{La revue de direction}

L'article~9.3 fait de la revue de direction le moment où la direction
regarde le SMSI dans son ensemble et décide. La norme en fixe les éléments
d'entrée~: l'état des actions décidées lors des revues précédentes~; les
changements des enjeux et des attentes des parties intéressées~; les
retours sur la performance (non-conformités et actions correctives,
résultats de la surveillance, des audits, atteinte des objectifs)~; les
retours des parties intéressées~; les résultats de l'appréciation des
risques et l'état du plan de traitement~; les opportunités
d'amélioration~\cite{iso27001_2022}.

\begin{boxinfo}
\textbf{La revue de direction de la BANQUE SANAGA.} Elle se tient le
30~mai 2028, présidée par M.~FOTSO. Le dossier comprend le tableau de bord
(Tableau~11.2) et le rapport de l'audit de mars. Il ne comprend pas le
rapport de l'audit interne du SMSI, remis le 2~juin. La revue décide des
moyens pour le mode dégradé et la sensibilisation, mais ne dit rien de
R-AG-02~: personne ne le lui a présenté.
\end{boxinfo}

% ============================================================================
\section{L'Audit de Certification et ses Constats}
\label{sec:c11-certification}
% ============================================================================

\subsection{Majeure, mineure}

\begin{boxdef}
\termecle{Non-conformité majeure}\index{non-conformité majeure}~:
non-conformité qui affecte la capacité du système de management à
atteindre les résultats escomptés~; par exemple une exigence qui n'est pas
appliquée du tout, ou un défaut répété, ou un doute sérieux sur la
maîtrise d'un risque~\cite{iso17021_2015}.

\medskip
\termecle{Non-conformité mineure}\index{non-conformité mineure}~:
non-conformité qui n'affecte pas cette capacité~: un défaut ponctuel, isolé,
dans un système qui fonctionne par ailleurs.

\medskip
\termecle{Piste de progrès}\index{piste de progrès}~: constat qui n'est pas
une non-conformité, mais une occasion d'améliorer.
\end{boxdef}

Les conséquences diffèrent. Une non-conformité \emph{mineure} donne lieu à
un plan d'actions que l'organisme accepte, et dont il vérifie l'effet au
prochain audit. Une non-conformité \emph{majeure} suspend la décision de
certification~: l'organisme attend la correction et la preuve de son
efficacité, dans le délai fixé par ses règles, et peut organiser un audit
complémentaire.

\subsection{Les constats de CERTIQUAL AFRIQUE}

\begin{tableaucours}{cL{5.6cm}cL{3.6cm}}{%
    \tableauHeader{N°} &
    \tableauHeader{Constat} &
    \tableauHeader{Exig.} &
    \tableauHeader{Qualification}
}
1 & Le risque des GAB de 14 agences, connu depuis avril, n'a pas été
apprécié ni traité & 6.1.2, 8.2 & \textbf{Majeure}~: exigence non appliquée
à une partie du périmètre, défaut déjà connu \\
2 & Mode dégradé non exercé de façon vérifiable dans 8 agences sur 38 &
Annexe~A, 5.30 & Mineure~: plan d'actions engagé \\
3 & La revue de direction n'a pas examiné les résultats de l'audit
interne & 9.3 & Mineure~: revue tenue, élément d'entrée manquant \\
4 & L'indicateur des correctifs ne couvre pas les postes des agences & 9.1
& Mineure \\
5 & Exercice de réponse aux incidents non réalisé depuis 2026 & Annexe~A,
5.24 & Piste de progrès~: planifié pour juillet \\
6 & La mesure 8.28 réintégrée n'a pas encore de procédure & Annexe~A, 8.28
& Piste de progrès \\
\end{tableaucours}
\vspace{-0.8em}
\begin{center}{\small\itshape Tableau~11.3. Les constats de l'étape~2
(juin 2028)}\end{center}

\begin{boxscenario}{Opération SANAGA. Pourquoi « majeure »}
\textbf{Analyse} (confiance élevée). Le constat~1 n'est pas grave par son
objet seul (quatorze agences, un équipement)~: il l'est parce qu'il montre
que le \emph{processus} d'appréciation des risques ne fonctionne pas comme
la norme l'exige. Un risque connu a attendu deux mois sans décision, et ni
l'audit interne, ni la revue de direction n'ont pu le faire remonter à
temps. C'est exactement le mécanisme de 2025~: le système laisse une zone
sans appréciation. Les constats~2 à~4, eux, sont des défauts dans un
système qui, par ailleurs, les a détectés.
\end{boxscenario}

\subsection{Ce qui suit}

La banque apprécie R-AG-02 en juillet, obtient du prestataire des GAB un
avenant de droit d'audit, et fait vérifier la correction par l'organisme.
CERTIQUAL AFRIQUE accepte la correction le 25~août et renouvelle le
certificat en septembre 2028. Entre la fin du certificat de 2025 et le
nouveau, la banque est restée trois mois sans certificat~: il a fallu
l'expliquer à deux correspondants étrangers et à un client institutionnel.

\begin{boxobjet}{R-AG-02 · Compromission d'un GAB (appréciée en juillet 2028)}
\textbf{Propriétaire~:} le directeur du réseau d'agences.

\textbf{Vraisemblance~:} 3 (système sans correctifs, exposé dans les
agences). Base~: configurations obtenues grâce à l'avenant (\OUI{}).

\textbf{Impact~:} 3 (retraits frauduleux plafonnés par les encaisses~;
indisponibilité). Base~: \OUI{}.

\textbf{Criticité~:} 9, intolérable~: remplacement du système des 14 GAB
programmé avant la fin de 2028~; en attendant, cloisonnement des GAB sur un
réseau dédié et surveillance des retraits atypiques.
\end{boxobjet}

\begin{boxinfo}
\textbf{La norme et la loi.} L'article~30 de la loi n°~2010/012 demande
aux exploitants d'évaluer et de réviser leurs systèmes de sécurité selon
l'évolution des technologies (Chapitre~\ref{chap:droit})~\cite{loi2010012}.
La revue de direction et l'audit interne du SMSI en sont la forme
normalisée~: leurs comptes rendus sont les pièces qu'un audit réglementaire
demandera.
\end{boxinfo}

% ============================================================================
\section{Ce que le SMSI demande à chaque Métier}
\label{sec:c11-metiers}
% ============================================================================

\begin{boxmetier}{Consultant en mise en œuvre de SMSI}
\textbf{Faire fonctionner, pas seulement documenter}
\begin{enumerate}
    \item Conduire la mise en œuvre dans l'ordre des étapes, en commençant
          par la décision de la direction et le périmètre.
    \item Rédiger un périmètre qui nomme ses exclusions et inclut ses
          interfaces.
    \item Mettre en place, dès le début, l'appréciation des risques en
          continu~: un risque nouveau est apprécié quand il apparaît, pas à
          la revue annuelle.
    \item Laisser au moins un cycle complet (audit interne, revue de
          direction, actions) avant l'audit de certification.
\end{enumerate}
\end{boxmetier}

\begin{boxmetier}{Auditeur en cybersécurité}
\textbf{Qualifier par le système, pas par l'objet}

Une non-conformité est majeure quand elle révèle que le système ne
fonctionne pas, pas quand son objet impressionne. L'auditeur justifie sa
qualification par l'exigence concernée et par ce que le défaut dit du
processus.
\end{boxmetier}

\begin{boxmetier}{RSSI}
\textbf{Un indicateur fiable vaut mieux qu'un bon indicateur}

La RSSI corrige I2 (inclure les postes), remplace la source de I5 par des
comptes rendus chronométrés, et présente I6 à la direction sans
l'enjoliver~: 11~\% de clics est un chiffre utile parce qu'il est vrai.
\end{boxmetier}

\begin{boxmetier}{Responsable GRC}
\textbf{Relier les constats d'audit au registre}

Tout constat d'audit qui révèle un risque nouveau entre dans le registre
dans la semaine, avec un propriétaire et une échéance d'appréciation. Le
défaut de la banque en 2028 tenait à un délai, pas à une ignorance.
\end{boxmetier}

\separateur

\begin{boxresume}
\textbf{Ce qu'il faut retenir de ce chapitre~:}
\begin{itemize}
    \item \textbf{Un SMSI} se met en œuvre comme un projet, puis vit comme
          un système dont les étapes se répètent.
    \item \textbf{Un périmètre défendable} nomme ses exclusions, inclut ses
          interfaces et est approuvé par l'instance qui répond des risques.
    \item \textbf{Les indicateurs} sont de mise en œuvre ou d'efficacité~;
          un indicateur vaut ce que vaut sa source.
    \item \textbf{L'audit interne} porte sur le système (articles~4 à~10) et
          sur les mesures~; la revue de direction décide à partir de ses
          résultats.
    \item \textbf{Majeure ou mineure}~: la qualification dépend de ce que
          le défaut dit de la capacité du système, non de son objet.
    \item \textbf{Une majeure} suspend la certification~: à la BANQUE
          SANAGA, trois mois sans certificat.
\end{itemize}
\end{boxresume}

\begin{boxinfo}
\textbf{Points de contrôle.} Vous avez compris ce chapitre si vous savez~:
\begin{itemize}
    \item rédiger un périmètre avec une exclusion défendable~;
    \item distinguer un indicateur de mise en œuvre d'un indicateur
          d'efficacité, et coter sa fiabilité~;
    \item composer l'ordre du jour d'une revue de direction~;
    \item qualifier un constat et justifier la qualification.
\end{itemize}
\end{boxinfo}

% ============================================================================
\section*{Exercices}
% ============================================================================

\begin{boxexo}[title={\ding{45}\quad Exercice~11.1. Exigé ou non~? \niveaubase}]
Dites si chacun des documents suivants est une information documentée
expressément exigée par l'ISO/IEC~27001, et citez l'article le cas échéant.
\begin{enumerate}[(a)]
    \item La déclaration d'applicabilité.
    \item Une charte d'usage de la messagerie.
    \item Le compte rendu de la revue de direction.
    \item Les attestations de formation des administrateurs.
    \item Un schéma de l'architecture réseau.
    \item La liste des non-conformités et des actions menées.
\end{enumerate}
\end{boxexo}

\begin{boxexo}[title={\ding{45}\quad Exercice~11.2. Construire un indicateur \niveaumoyen}]
\begin{enumerate}[(a)]
    \item Classez en mise en œuvre ou efficacité~: part des agences dont le
          partage est cloisonné~; nombre de tentatives d'accès entre agences
          bloquées par mois~; délai moyen de détection d'un incident~; part
          des contrats de prestataires comportant une clause de sécurité.
    \item Rédigez la fiche complète d'un indicateur d'efficacité pour la
          mesure 5.30 (mode dégradé), avec sa source et sa cote de
          fiabilité.
    \item Corrigez la définition ou la source de I2 pour qu'il devienne
          \OUI{}.
\end{enumerate}
\end{boxexo}

\begin{boxexo}[title={\ding{45}\quad Exercice~11.3. Qualifier \niveaumoyen}]
Qualifiez chacun des constats suivants (majeure, mineure, piste de
progrès) et justifiez par ce qu'il dit du système.
\begin{enumerate}[(a)]
    \item Aucun audit interne du SMSI n'a été conduit depuis la
          certification initiale.
    \item Sur 25 comptes de partants prélevés, un est resté actif huit jours
          au lieu de cinq~; les 24 autres sont conformes.
    \item La politique de sécurité n'a pas été revue depuis trois ans,
          alors que la procédure prévoit une revue annuelle~; elle reste
          pertinente.
    \item Les objectifs de sécurité de 2028 ne sont pas mesurables
          («~renforcer la sécurité~»).
    \item Le tableau de bord pourrait comporter un indicateur sur les
          fournisseurs.
\end{enumerate}
\end{boxexo}

\begin{boxexo}[title={\ding{45}\quad Exercice~11.4. Le SMSI de LiZbiZ \niveauexpert}]
LiZbiZ décide d'engager un SMSI en vue d'une certification fin 2029
(Chapitre~\ref{chap:iso27000}, Exercice~8.4).
\begin{enumerate}[(a)]
    \item Établissez le calendrier des huit étapes du Tableau~11.1, en
          laissant un cycle complet avant l'audit initial.
    \item Proposez quatre indicateurs (deux de mise en œuvre, deux
          d'efficacité) qui couvrent le scénario SO1, avec leur source.
    \item LiZbiZ n'a pas d'audit interne. Qui peut conduire l'audit interne
          du SMSI, et à quelles conditions d'indépendance
          (Chapitre~\ref{chap:droit})~?
    \item Rédigez l'ordre du jour de la première revue de direction, en
          couvrant tous les éléments d'entrée exigés.
\end{enumerate}
\end{boxexo}

\bigskip

\begin{boxinfo}
\textbf{Pour aller plus loin.}
\begin{itemize}
    \item L'ISO/IEC~27001:2022~\cite{iso27001_2022}, articles~4.3, 9 et~10.
    \item L'ISO/IEC~27004:2016~\cite{iso27004_2016}, sur les mesures.
    \item L'ISO/IEC~27007:2020~\cite{iso27007_2020}, sur l'audit d'un SMSI.
    \item L'ISO/IEC~17021-1:2015~\cite{iso17021_2015}, sur les exigences
          applicables aux organismes de certification.
\end{itemize}
\end{boxinfo}

\begin{boxinfo}
\textbf{Transition.} L'audit d'un SMSI vérifie le système et ses mesures,
souvent à partir de documents et d'extractions. Il faut parfois aller plus
loin, jusqu'à la configuration des serveurs, la segmentation du réseau,
les comptes et les vulnérabilités. C'est l'audit technique, objet du
Chapitre~\ref{chap:audit-technique}.
\end{boxinfo}
